\chapter{Annexos}
\label{ch:annexos}

En aquest apartat s'inclouen materials addicionals que complementen la memòria del projecte.

\section{Exemple d'esquema de configuració (Schema-driven JSON)}
\label{annex:schema_json}
A continuació es mostra un exemple simplificat de l'estructura JSON que el \textit{backend} (\textbf{la-core-legacy}) envia al \textit{frontend} per construir dinàmicament la vista de configuració d'una càmera. S'observa com es defineix l'estructura, el tipus de dada, els rangs de valors i la persistència que s'ha d'aplicar.

\begin{lstlisting}[language=JavaScript, caption={Exemple d'esquema JSON per a la configuració d'una càmera IP.}, label={lst:schema_json}]
{
  "properties": {
    "camera": {
      "type": "object",
      "properties": {
        "id": {
          "type": "integer",
          "readOnly": true
        },
        "name": {
          "type": "string",
          "maxLength": 64,
          "description": "Nom assignat a la càmera"
        },
        "enabled": {
          "type": "boolean",
          "default": true
        },
        "stream": {
          "type": "object",
          "properties": {
            "codec": {
              "type": "string",
              "enum": ["H264", "H265"],
              "default": "H264"
            },
            "resolution": {
              "type": "string",
              "enum": ["1080p", "720p", "4K"]
            },
            "fps": {
              "type": "integer",
              "minimum": 1,
              "maximum": 60,
              "default": 25
            }
          }
        }
      }
    }
  }
}
\end{lstlisting}

Aquest esquema permet que el backend defineixi regles de validació (com el regex d'IP) que el frontend aplica automàticament sense conèixer prèviament el camp.
